In this project, we developed two core pipelines for SpurSol's Knoccs CRM platform: a Design-to-Code Pipeline and a RAG-based Knowledge Base. We detailed our methodology for iterative prompt engineering across nine instruction file iterations, multi-agent orchestration with eight specialized agents, component generation and evaluation, Figma-to-code synchronization, and RAG system design with hybrid retrieval and DeepEval-based evaluation.
\\\\
The results demonstrate the feasibility of AI-assisted design-to-code automation in a production Angular codebase: 20 components were generated at an average CLIP visual similarity of 82.6\% (computed per component, comparing one Storybook screenshot against its corresponding Figma design). The RAG-based knowledge retrieval system was evaluated across 52 ground truth queries spanning 10 companies, and integrated into the Knoccs platform. Tag and Macro Management features were implemented to validate that the generated component library could support real feature development within the production application.
\\\\
While the progress thus far is promising, several avenues remain for exploration. We intend to restructure the Knoccs design system with machine-readable semantics to reduce token usage and improve generation accuracy, and to benchmark alternative LLMs (GPT-4o, Gemini, open-source models) against Claude for component generation quality and cost. For the RAG pipeline, connecting to Knoccs's live production data through an automated ingestion pipeline is the most impactful next step, alongside migrating to gRPC communication to align with the existing microservices architecture. Further improvements to the embedding model, chunking strategy, and LLM selection are planned to enhance retrieval precision and response quality. Continuous evaluation against expanding ground truth datasets and rigorous validation within real client workflows will be central to the next stages of development.
